\section{The finite-size engine: statement and setting}\label{sec:statement}
Fix a finite local dimension $q$ and three Hermitian matrices $h_0,h_h,h_v$,
of sizes $q$, $q^2$, $q^2$, respectively. On
$\Lambda_L=\{0,\ldots,L-1\}^2$, let
\begin{equation}\label{eq:model}
 H_L=\sum_{x\in\Lambda_L}(h_0)_x+
 \sum_{\substack{x,x+e_1\in\Lambda_L}}(h_h)_{x,x+e_1}+
 \sum_{\substack{x,x+e_2\in\Lambda_L}}(h_v)_{x,x+e_2}.
\end{equation}
The orientation of each bond is fixed; rotation invariance is not required.
There are no additional boundary operators. When the spectrum has at least
two distinct eigenvalues, set
\begin{equation}\label{eq:gap}
 E_0(L)=\min\spec H_L,\qquad
 \Delta(L)=\min\{\lambda-E_0(L):\lambda\in\spec H_L,\ \lambda>E_0(L)\}.
\end{equation}
We assume this gap is defined for all sufficiently large $L$. Thus scalar
finite-volume Hamiltonians are excluded when necessary, and a degenerate
ground space does not itself make the gap zero. For $L\geq2$ put
\begin{equation}\label{eq:exponents}
 z_L=\frac{-\log\Delta(L)}{\log L},\qquad
 z^-=\liminf_{L\to\infty}z_L,\qquad
 z^+=\limsup_{L\to\infty}z_L.
\end{equation}
An existing finite exponent is $z=z^-=z^+$.

For a real $a$, the conventions $a\in\LS$ and $a\in\LP$ refer to
the arithmetical complexity of the strict lower rational cut
$\{r\in\Q:r<a\}$. A real is in $\Dtwo$ if it is computable relative
to the halting oracle, equivalently if it is a limit of a computable
rational sequence. These are representations of real numbers, not
conventions about an unspecified digit expansion
\cite{zhengWeihrauch,density}.

\begin{theorem}[Exact classification]\label{thm:classification}
Consider \eqref{eq:model} with integer interaction matrices and arbitrary
finite local dimension.
\begin{enumerate}
\item The existing exponents in $[0,\infty)$ are exactly
      $\Dtwo\cap[0,\infty)$.
\item The finite pairs $(z^-,z^+)$ in $[0,\infty)^2$ are exactly
      the pairs $(a,b)$ with $0\leq a\leq b<\infty$,
      $a\in\LS$ and $b\in\LP$.
\end{enumerate}
Each pair is realized by one fixed Hamiltonian. Its local dimension may
depend on the pair. The necessity statements hold more generally for
algebraic interaction matrices, in any fixed dimension, whenever the gaps
and endpoints in question are defined and finite.
\end{theorem}

\begin{corollary}[One fixed uncomputable exponent]\label{cor:omega}
For any chosen prefix-universal machine $U$, there is a fixed model
\eqref{eq:model}, with integer matrices, for which
$\Delta(L)=L^{-\Omega_U+o(1)}$, where $\Omega_U$ is the halting probability
of $U$. No uncomputable number occurs in a matrix entry.
\end{corollary}

The construction and its full spectral estimate are in
Sections~\ref{sec:classical}--\ref{sec:assembly}. Section~\ref{sec:words}
provides the required computationally fast words, and
Section~\ref{sec:necessity} proves the converse.

\section{Two elementary spectral and boundary estimates}\label{sec:lemmas}
We first record estimates that will control all computational and
noncomputational sectors. Graph Laplacians use the quadratic form
$\sum_{uv\in E}|\psi_u-\psi_v|^2$, with repeated edges counted with their
integer weights.

\begin{lemma}[Penalty localization]\label{lem:penalty}
Let $C$ be a finite connected graph of maximum degree at most $D$, and let
$V$ be diagonal with values in $\{0\}\cup[1,\infty)$. If all vertices
are penalized, then $L_C+V\geq I$. If both free and penalized vertices
occur, and each connected free cluster has at most $g$ vertices, then
\begin{equation}\label{eq:penalty}
 L_C+V\ \geq\ \frac{1}{8Dg^2}I.
\end{equation}
A free path on $n\geq2$ vertices has positive gap
$2-2\cos(\pi/n)\geq4/n^2$; a free cycle on $n\geq3$ vertices has positive
gap $2-2\cos(2\pi/n)\geq4/n^2$. A free isolated vertex has no positive
eigenvalue. The two-vertex cycle with two parallel unit edges has gap $4$.
\end{lemma}
\begin{proof}
The all-penalized case is immediate. For each free cluster $K_i$, choose an
edge $u_ip_i$ to a penalized vertex. Each penalized vertex is selected at
most $D$ times, so half of its penalty can allocate $1/(2D)$ to each
selected cluster. The remaining half pays for its own norm. Therefore the
energy of $\psi$ is at least
\[
 \sum_i\left(\sum_{uv\in E(K_i)}|\psi_u-\psi_v|^2+
 |\psi_{u_i}-\psi_{p_i}|^2+\frac{|\psi_{p_i}|^2}{2D}\right)
 +\frac12\|\psi_{C\setminus F}\|^2,
\]
where $F$ is the free set. Let $M_i=\max_{v\in K_i}|\psi_v|$.
If $|\psi_{p_i}|\geq M_i/2$, the allocated penalty is at least
$M_i^2/(8D)$. Otherwise a path from a maximum to $p_i$, of at most $g$
edges, has energy at least $M_i^2/(4g)$ by Cauchy--Schwarz. Since
$\|\psi_{K_i}\|^2\leq gM_i^2$, either case bounds the cluster norm by
\eqref{eq:penalty}. The path and cycle statements follow from their
explicit spectra. Keeping parallel-edge weights gives the last case;
larger integer weights cannot reduce the lower bounds.
\end{proof}

\begin{lemma}[Boundary preference]\label{lem:boundary}
At each site put four flag bits $T,F,B,R$. Give a site energy
$-(T+F+B+R)$. On a vertical bond, with upper endpoint $x$ and lower
endpoint $y$, add $2T_y+2B_x$; on a horizontal bond with left endpoint
$x$ and right endpoint $y$, add $2F_y+2R_x$. Add unit penalties for
forbidden onsite and bond patterns. Then
\begin{equation}\label{eq:flags}
 E_{\mathrm{cl}}(\sigma)+4L
 =\#\{\text{incorrect flag bits}\}+\#\{\text{violated local rules}\}.
\end{equation}
Here a correct flag is the indicator of its corresponding physical edge.
If exactly one configuration has correct flags and no rule violation, it
is the unique classical minimizer, of energy $-4L$, and every other
configuration has energy at least $-4L+1$.
\end{lemma}
\begin{proof}
The coefficient of $T$ is $-1$ on the top row and $+1$ elsewhere,
because every other site has one bond from above. Its energy is therefore
$-L$ plus the number of errors relative to the top-row indicator.
Apply the same argument separately to $F,B,R$ and add the rule penalties.
\end{proof}
All bond contributions in this lemma are genuine two-site matrices, even
when one factor is the identity. The edge preference comes solely from
missing bonds in \eqref{eq:model}. We have proved a classical gap
\emph{at least} one.

\section{A unique classical tableau from uniform interactions}\label{sec:classical}
Fix an integer $\beta\geq1$ and a computable binary word $A(0),A(1),\ldots$
whose bits can be written sufficiently quickly; Section~\ref{sec:words}
constructs every word we need. Fix a total deterministic one-tape machine
$M$ that writes the word on a persistent output track, with bit $j$ at
column $j$, and never crosses the left edge. It can be arranged that the
first $j$ bits are written by time $2^{j-1}-2$ for every sufficiently large
$j$. This is a timing upper bound, not a time supplied as a coupling.

Besides the four flags, the classical alphabet contains a binary digit
$b$, a carry $c$, a past-one flag $f$, and the tape, head, persistent output
and arrival fields of $M$. Every listed rule is enforced by a diagonal
onsite or nearest-neighbour unit penalty. Conjunctions can be imposed as
a single forbidden-pattern projector; their precise grouping does not
affect \eqref{eq:flags}.

\subsection{Counting the row number}
The counter rules are
\begin{align*}
 T=1&\ \Longrightarrow\ b=f=0,&
 F=1&\ \Longrightarrow\ c=1,\\
 F_y=0&\ \Longrightarrow\ c_y=b_xc_x
      &&\text{on a horizontal bond }x\to y,\\
 T_y=0&\ \Longrightarrow\ b_y=b_x\mathbin{\mathrm{XOR}}c_x,
 \quad f_y=f_x\mathbin{\mathrm{OR}}b_x
      &&\text{on a vertical bond }x\to y.
\end{align*}
With correct flags, carries are fixed left to right on each row, and the
vertical rule advances the binary row number. Induction gives row $r$
equal to the binary expansion of $r$, with least significant bit at column
zero; $f$ records whether that column was one on an earlier row. Thus the
derived label
\begin{equation}\label{eq:register}
 \reg=B\,(f\mathbin{\mathrm{OR}}b)
\end{equation}
selects exactly columns $0,\ldots,\ell(L)-1$ on the bottom row, where
\begin{equation}\label{eq:ell}
 \ell(L)=\bitlength(L-1)=\log_2 L+O(1).
\end{equation}
Indeed a column $j$ has shown a one by row $L-1$ exactly when
$2^j\leq L-1$. For $L\geq2$, $\ell(L)<L$.

\subsection{A nearest-neighbour computation tableau}
Let the tape field be $\gamma$, the head state be $h\in Q\cup\{\blank\}$,
the output be $o\in\{\blank,0,1\}$, and the two arrival signals be
$a_L,a_R\in Q\cup\{\blank\}$. The top row has blank tape and output,
with the initial head only at its leftmost site. On each horizontal bond,
the left site's transition fixes the right site's $a_L$; the right site's
transition fixes the left site's $a_R$. A missing source gives a blank
arrival. Require $a_L=\blank$ at $F=1$ and $a_R=\blank$ at $R=1$.
These equalities fix every unused arrival signal as well.

On a vertical bond, the upper site's own transition fixes the tape and
output below it. Its staying-head transition, or its $a_L$ or $a_R$,
fixes the next head; two incoming heads are forbidden. This is a finite
rule table determined by the transition function of $M$. In particular,
the arrival fields implement the update by horizontal and vertical
\emph{two-site} tests; a three-site update is not an interaction.

Given one row's tape, head and output, first fix all arrival fields by the
horizontal tests and edge initialization, then fix the next row by the
vertical tests. This proves uniqueness by induction, including the last
row's arrival fields. At time $t$ the head column is at most $t$.
A transition to another included row cannot leave the right side of an
$L\times L$ tableau. An outgoing move after the final row places no
condition on a nonexistent next row. Consequently the tableau exists and
is unique for every $L\geq2$, regardless of whether all desired output
bits have yet appeared.

Together with the counter, this gives a unique correct classical
configuration $\sigma_L^*$ and the unit separation of
Lemma~\ref{lem:boundary}. On the bottom register put
\begin{equation}\label{eq:radix}
 m_j=\beta\ind_{\{o_j=1\}},\qquad r_j=2^{m_j},\qquad 0\leq j<\ell(L).
\end{equation}
For all sufficiently large $L$, the timing assumption gives $o_j=A(j)$
for $j<\ell(L)$: the first $\ell$ bits appear by row
$2^{\ell-1}-2$, and $L-1\geq2^{\ell-1}$. Therefore
\begin{equation}\label{eq:k}
 k(L)=\sum_{j<\ell(L)}m_j
     =\beta\,|A\cap[0,\ell(L))|.
\end{equation}

\section{The quantum counter and all its sectors}\label{sec:counter}
At every site use the same additional space with digit
$d\in\{0,\ldots,2^\beta-1\}$ and head-side symbol
$s\in\{\NL,\NR,\INC,\RET\}$. The symbols $\NL$ and $\NR$ mean
left and right of a head; the other two are head states. Outside $\reg$,
penalize every dynamic field other than $(0,\NL)$ by one and apply no
transition. Inside $\reg$, penalize digits $d\geq r_j$ and, on two
consecutive register sites, the side-symbol patterns
\[
 (\NR,\NL),\qquad (\text{head},\NL),\qquad
 (\NR,\text{head}),\qquad (\text{head},\text{head}).
\]
Either head state is meant here. All these penalties are nonnegative.

For each transition $u\leftrightarrow v$ below, add its edge Laplacian
$(|u\rangle-|v\rangle)(\langle u|-\langle v|)$. Digits not explicitly
changed are retained. The rules on the intended register are
\begin{align}
 (d,\INC)&\leftrightarrow(d+1,\RET)
              &&(d+1<r_j),\label{eq:increment}\\
 (r_j-1,\INC)(e,\NR)&\leftrightarrow(0,\NL)(e,\INC)
              &&(j+1<\ell),\label{eq:carry}\\
 (r_j-1,\INC)&\leftrightarrow(0,\RET)
              &&(j=\ell-1),\label{eq:wrap}\\
 (0,\NL)(e,\RET)&\leftrightarrow(0,\RET)(e,\NR)
              &&(j+1<\ell),\label{eq:return}\\
 (d,\RET)&\leftrightarrow(d,\INC)
              &&(j=0).\label{eq:restart}
\end{align}
The conditions involving $j$ in this notation are implemented as follows.
Rules~\eqref{eq:increment} and \eqref{eq:restart} are onsite, guarded by
$\reg=1$, with the latter also requiring $F=1$.
Rules~\eqref{eq:carry} and \eqref{eq:return} require both endpoints to be
register sites. Rule~\eqref{eq:wrap} is a horizontal bond term: its left
endpoint has $\reg=1$ and its right endpoint has $\reg=0$; it changes
only the left dynamic field and is the identity on the right dynamic
field. Such a right neighbour exists on $\sigma_L^*$ since $\ell(L)<L$.
Every digit test uses finite local classical labels. Thus every rule is
onsite or nearest-neighbour, including the last-cell test. No rule queries
coordinates or the absence of an unrepresented neighbour.

Write $K_L$ for this nonnegative dynamic operator on the correct classical
block. In the next lemma set $B_* =\prod_{j<\ell}r_j=2^k$; this symbol
is unrelated to the bottom flag $B$.

\begin{lemma}[Full counter gap]\label{lem:counter}
For $\ell\geq1$ the dynamic kernel is nonempty, and the smallest positive
eigenvalue obeys
\begin{equation}\label{eq:countergap}
 \frac{1}{64\ell^2 B_*^2}\ \leq\ \Delta^+(K_L)
 \ \leq\ \frac{4\pi^2}{B_*^2}.
\end{equation}
There is a free counter orbit of length
\begin{equation}\label{eq:orbit}
 N=2\sum_{j=0}^{\ell-1}\frac{B_*}{r_0\cdots r_{j-1}},
 \qquad 2B_*\leq N\leq2\ell B_*.
\end{equation}
The empty product is one. If $N=2$, the orbit retains two parallel edges.
\end{lemma}
\begin{proof}
First fix the exterior fields. A wrong exterior field cannot change under
any transition and carries a penalty at least one, so the entire
corresponding block is at least one. With the exterior correct, head number
on the register is conserved.

With zero heads there are no transitions. Free isolated configurations,
including the all-zero, all-$\NL$ configuration, give ground states;
penalized isolated configurations cost at least one. With two or more
heads every configuration is penalized: between consecutive heads the
side-symbol rules would require $\NR$ immediately after the first and
$\NL$ immediately before the second, forcing a forbidden switch.
Adjacent heads are forbidden directly. The all-penalized part of
Lemma~\ref{lem:penalty} applies even if this sector has large degree.

With one head, there is at most one forward match and one backward match.
At an incrementing head the forward choice is increment, carry or wrap;
its predecessor is a carry or restart. At a returning head, its digit
distinguishes an increment predecessor (nonzero) from a wrap or return
predecessor (zero, at or away from the last site). The return or restart
successor is also unique. Illegal digits or side strings remove possible
matches. Hence transition degree is at most two, counting weights.

There are exactly $2\ell B_*$ free one-head configurations: choose a head
position, one of its two states, and legal digits; the side symbols are
forced. On a component meeting penalties, a free cluster has size at most
$2\ell B_*$, so Lemma~\ref{lem:penalty}, with $D=2$, gives
$1/[16(2\ell B_*)^2]$. On a completely free component, a nontrivial
path or cycle has at most $2\ell B_*$ vertices and the positive spectral
gap is larger than the same bound. Free singletons only enlarge the
ground space. These arguments prove the lower bound in every sector.

For the upper bound start with an incrementing head at site zero and any
legal digit string. It increments the mixed-radix value, carries over
maximal digits, and returns over the zeros it has created. Overflow wraps
to zero. It returns to the initial digit value after $B_*$ increments.
Each cell $j$ is visited in the forward sweep $B_*/(r_0\cdots r_{j-1})$
times and contributes twice that many transitions in a complete orbit.
This gives \eqref{eq:orbit}. The orbit is free and its positive gap is
$2-2\cos(2\pi/N)\leq4\pi^2/B_*^2$.
When $\ell=1,r_0=1$, wrap and restart contribute to the same pair of
states, giving the double-edge cycle with gap four. This case satisfies
the same inequalities.
\end{proof}

We use the orbit only for an upper bound. Equality of the full gap with
its cycle eigenvalue is unnecessary and is not asserted. In particular,
other free one-head components and their ground states are allowed.

\section{Assembly and the realization formula}\label{sec:assembly}
The complete local space is the tensor product of the finite classical
alphabet and the dynamic alphabet. Add the classical terms to the
guarded dynamic terms. All matrix entries are integers: diagonal rules
are integer projectors and each transition contributes diagonal $+1$
and off-diagonal $-1$. The matrices are fixed once $M$ and $\beta$ are
chosen. They have exactly the form \eqref{eq:model}.

\begin{theorem}[Boundary realization]\label{thm:realization}
For every word satisfying the timing condition of
Section~\ref{sec:classical} and every integer $\beta\geq1$, the constructed
Hamiltonian has, for all sufficiently large $L$,
\begin{equation}\label{eq:gapbracket}
 \frac{4^{-k(L)}}{64\ell(L)^2}\ \leq\ \Delta(L)
 \ \leq\ \min\{1,4\pi^2\,4^{-k(L)}\}.
\end{equation}
Consequently, writing $\rho_J(A)=|A\cap[0,J)|/J$,
\begin{equation}\label{eq:densityformula}
 z_L=2\beta\rho_{\ell(L)}(A)
      +O\!\left(\frac{\log\log L}{\log L}\right).
\end{equation}
In particular $z^-=2\beta\liminf_J\rho_J(A)$ and
$z^+=2\beta\limsup_J\rho_J(A)$.
\end{theorem}
\begin{proof}
All classical labels are conserved. The exact block decomposition is
\begin{equation}\label{eq:blocks}
 H_L=\bigoplus_\sigma\bigl(E_{\mathrm{cl}}(\sigma)I+K_\sigma\bigr),
 \qquad K_\sigma\geq0.
\end{equation}
The correct block has energy $-4L$ and a nonempty kernel. Every other
classical block is at least $-4L+1$, independently of its apparent
register, head pattern or output. A superposition across blocks cannot
evade this operator inequality. Within the correct block exterior defects,
zero-head excitations, multi-head sectors and all one-head components were
covered in Lemma~\ref{lem:counter}. Thus
\[
 \min\{1,\Delta^+(K_L)\}\leq\Delta(L)\leq\Delta^+(K_L).
\]
The lower estimate in \eqref{eq:countergap} is at most one and proves the
lower bound in \eqref{eq:gapbracket}. Changing a single frozen exterior
field from $(0,\NL)$ to $(0,\NR)$ in a dynamic ground state produces a
unit-energy eigenvector. There are no register/exterior adjacency
penalties and the wrap term is the identity on the exterior dynamic
field, so this gives $\Delta(L)\leq1$. The orbit gives the other upper
bound.

Taking base-two logarithms of the bracket gives
$\log_2(1/\Delta(L))=2k(L)+O(\log\ell(L))$. Equations
\eqref{eq:ell} and \eqref{eq:k} yield \eqref{eq:densityformula}.
Every sufficiently large integer occurs as $\ell(L)$, and $\ell(L)$
increases only by unit steps, which proves the endpoint equalities.
\end{proof}

Wrong tableaux, fake interior boundaries and arbitrarily long apparent
computations are covered by the classical unit separation, without any
runtime bound on those wrong blocks. The only true computation is static.
The closing eigenvectors keep that tableau fixed, freeze the exterior,
and evolve the quantum counter on the logarithmic boundary segment.

\section{Fast words with prescribed density endpoints}\label{sec:words}
The usual computable-density classification supplies the endpoint
arithmetic \cite{density}. The following timing refinement is needed
because only $L$ rows are available to write approximately $\log_2 L$
output bits.

\begin{lemma}[Fast density word]\label{lem:word}
If $0\leq a\leq b\leq1$, $a\in\LS$ and $b\in\LP$, there is a binary
word $A$ with lower density $a$ and upper density $b$ whose $j$th bit is
computable in time polynomial in $\log(j+2)$. The polynomial and its
constants may depend on the endpoints' programs.
\end{lemma}
\begin{proof}
The sequence characterizations in
\cite[Theorem~5.1 and Corollaries~5.4--5.5]{density} give computable
rational sequences with liminf $a$ and limsup $b$. If $a<b$, choose a
fixed rational $c$ with $a<c<b$. Clip the first sequence to $[0,c]$ and
the second to $[c,1]$, and interleave them. Continuity and monotonicity
of clipping preserve the desired endpoint of each sequence. The result
is a computable rational sequence $a_n\in[0,1]$ with endpoints $a,b$.
If $a=b$, membership of both cuts gives a $\Dtwo$ real, so use a
convergent computable rational sequence clipped to $[0,1]$. Choosing the
separator uniformly from arbitrary endpoint names is not claimed.

Fix $a_0=0$. Put $\sigma(0)=0$. At stage $s$, simulate the canonical
algorithm for $a_{\sigma(s)+1}$ from its start for $s$ steps, and advance
$\sigma$ by one exactly when the output has appeared. Otherwise retain
its value. Totality makes $\sigma$ unbounded, and unit advances preserve
liminf and limsup for $b_s=a_{\sigma(s)}$. Every selected output other
than $a_0$ has $O(s)$ bits and comes from a bounded simulation.
Reconstructing all advances to stage $s$ takes $O(s^2)$ simulated steps,
with polynomial universal-simulation overhead. This explicit output
bound prevents a large rational name from defeating the slowdown.

Let $N_0=1$ and $N_{s+1}=(s+2)N_s$. Set $A(0)=0$ and, for
$0\leq t<N_{s+1}-N_s$, put
\begin{equation}\label{eq:word}
 A(N_s+t)=\lfloor(t+1)b_s\rfloor-\lfloor t b_s\rfloor.
\end{equation}
The number of ones in the first $t$ letters of that stage is
$\lfloor t b_s\rfloor$. If $J=N_s+t$, its density differs by less than
$1/J$ from
\begin{equation}\label{eq:convex}
 \frac{N_s}{J}\rho_{N_s}(A)+\frac{t}{J}b_s.
\end{equation}
At the end of the stage, the density differs from $b_s$ by at most
$1/(s+2)+1/N_{s+1}$. Endpoints therefore attain the liminf and limsup
of $b_s$. The convex formula shows that intermediate densities cannot
create additional limit extrema. This proves both density assertions,
including $a=0$ or $b=1$.

To compute $A(j)$, find its factorial-sized stage, reconstruct
$\sigma(s)$ and its bounded rational output, then evaluate the two floor
values in \eqref{eq:word}. The stage has $s=O(\log(j+2))$, all index
arithmetic has $O(\log(j+2))$ bits, and the selected rational has
polynomially bounded length. All steps take polynomial time in
$\log(j+2)$.
\end{proof}

\begin{proof}[Realization in Theorem~\ref{thm:classification}]
Given $(z^-,z^+)=(a,b)$ as in the theorem, choose an integer
$\beta\geq\max\{1,b/2\}$. Positive rational scaling preserves the cut
classes. Apply Lemma~\ref{lem:word} to $(a/(2\beta),b/(2\beta))$.
A one-tape machine can compute the bits and write them successively on
a persistent track, using separate work symbols. Simulation and travel
to column $j$ cost at most polynomial time in $j$ per bit. Its prefix
runtime is therefore eventually less than $2^{j-1}-2$.
Theorem~\ref{thm:realization} gives the prescribed pair. Taking $a=b$
gives every existing nonnegative $\Dtwo$ exponent.
\end{proof}

\begin{proof}[Proof of Corollary~\ref{cor:omega}]
Dovetail the prefix-universal machine's programs and enumerate its halting
weights, giving computable increasing rational sums $\Omega_{U,s}$
converging to $\Omega_U$. Use the convergent sequence
$\Omega_{U,s}/2$ in the word construction and set $\beta=1$.
The resulting exponent is $2(\Omega_U/2)=\Omega_U$. The finite simulator
and word algorithm determine the finite classical alphabet and integer
matrices. The real is encoded in an asymptotic density, not as a local
coefficient.
\end{proof}

An exponent zero does not assert a uniform positive gap. For example,
the all-radix-one counter has orbit length $2\ell(L)$ and a closing
gap of order $1/\ell(L)^2$, still with exponent zero.

\section{Necessity from finite algebraic spectra}\label{sec:necessity}
\begin{proposition}\label{prop:necessity}
For fixed algebraic local matrices, every defined finite existing exponent
is $\Dtwo$. If both endpoints are finite, the lower endpoint is in
$\LS$ and the upper endpoint is in $\LP$. This holds in any fixed
lattice dimension and for any uniformly computable choice of finite
boxes.
\end{proposition}
\begin{proof}
For each $L$, construct $H_L$ exactly over its algebraic coefficient
field. Exact characteristic polynomials and real-root isolation compute
its distinct eigenvalues and decide multiplicities. Thus the gap above
the entire ground space is a uniformly computable positive algebraic
number whenever defined. The hypothesis that it is eventually defined
is essential; the zero Hamiltonian has no second distinct eigenvalue.
We do not replace algebraic entries by arbitrary computable real entries
while retaining a multiplicity algorithm.

Compute rational $q_L$ within $1/L$ of $z_L$. This preserves all finite
endpoints. For rational $r$,
\[
 r<\liminf_Lq_L\quad\Longleftrightarrow\quad
 \exists t\in\Q\ \exists N\quad
 (r<t\ \text{and}\ \forall L\geq N:\ q_L>t).
\]
This is a $\Sigma^0_2$ lower cut. Dually the strict upper cut of
$\limsup_Lq_L$ is $\Sigma^0_2$. If that endpoint is irrational, its
strict lower rational cut is the complement of its strict upper cut,
hence $\Pi^0_2$; if it is rational, its lower cut is computable and is
also $\Pi^0_2$. This avoids a false complement identity at a rational
endpoint. An existing limit is computable in the halting oracle: that
oracle can find arbitrarily accurate tail-Cauchy intervals of the
computable sequence. Equivalently it is a $\Dtwo$ real.

For a $d$-dimensional box the norm estimate
$\Delta(L)\leq2\|H_L\|\leq C L^d$ gives an eventual lower bound for
$z_L$. A finite limsup gives bounded tails when needed. No efficient
diagonalization or computable rate of asymptotic convergence is claimed.
\end{proof}

This proves the remaining direction of
Theorem~\ref{thm:classification}. More generally, the same argument applies
to total computable postprocessing of uniformly computable finite-scale
spectral data. Any logarithm argument must be positive and every scale
denominator must be nonzero. Merely naming a scale does not supply these
domain conditions.

